\documentclass{article}
\usepackage{hyperref}
\usepackage{Style}

\nocite{*} % Comentar si quiero citar
%\addbibresource{bibliografia.bib} % Quitar el comentado si quiero usar bibliografia

\begin{document}

\begin{minipage}{2.5cm}
    \includegraphics[width=2cm]{imagen_puc.jpg}
\end{minipage}
\begin{minipage}{14cm}
    {\sc Pontificia Universidad Católica de Chile\\
    Facultad de Matemáticas\\
    Departamento de Matemática\\
    Profesor: Manuel Cabezas -- Ayudante: Benjamín Mateluna}
\end{minipage}
\vspace{1ex}

{\centerline{\bf Introducción a la Geometría - MAT1304}
\centerline{\bf Ayudantía 2}}
\centerline{\bf 12 de marzo de 2026}

\vspace{1cm}
\noindent\textbf{Problema 1.} Demuestre que dado $\triangle ABC$, entonces es isósceles de base 
$\overline{BC}$ si y solo si $\angle ABC=\angle ACB$. Deduzca, de la demostración previa, que la 
bisectriz, la altura y la mediana desde el punto opuesto a la base son iguales en un triángulo 
isósceles.

\vspace{5mm}
\begin{dfn}
    Dado un segmento $\overline{AB}$ y $M$ su punto medio, se dice que $L$ es simetral si y solo 
    si es perpendicular a $\overline{AB}$ y pasa por el punto $M$.
\end{dfn}

\vspace{2mm}
\noindent\textbf{Problema 2.} Sea $\overline{AB}$ un segmento y $L$ su simetral, demuestre que 
dado $P\in L$, se tiene que $\overline{AP}=\overline{BP}$.

\vspace{2mm}
\noindent\textit{Observación:} Notar que $L$ es el conjunto de puntos tales que $\triangle ABP$ es
isósceles.

\vspace{5mm}
\noindent\textbf{Problema 3.} Dado $\triangle ABC$, denotamos por $M$ al punto medio de 
$\overline{AB}$. Demuestre que si $\overline{AM}=\overline{BM}=\overline{CM}$, entonces 
$\angle BCA=90^{\circ}$.

\vspace{5mm}
\noindent\textbf{Problema 4.} Sea $\triangle ABC$ un triángulo equilátero, sea $a$ la medida del
lado del triángulo. Calcule el área del triángulo.

%\printbibliography % Quitar el comentado si quiero usar bibliografia

\end{document}
